\documentclass[11pt]{letter}
\usepackage[margin=1in]{geometry}
\usepackage{xcolor}
\definecolor{IEEEblue}{HTML}{00629B}
\usepackage[colorlinks=true,urlcolor=IEEEblue]{hyperref}
\signature{Yassir AL-Karawi\\Hamed Al-Raweshidy (Corresponding Author)}
\address{College of Engineering, University of Diyala, Iraq\\
Wireless Networks and Communications Centre, Brunel University of London, U.K.}
\date{13 July 2026}

\begin{document}
\begin{letter}{Editorial Office\\
Publication venue to be selected before submission}

\opening{Dear Editor,}

Please consider the following manuscript for publication:

\begin{center}
\emph{KPM-Bridge: An Uncertainty-Aware Cross-Implementation\\
Telemetry Fabric for Portable AI xApps in Open 6G RAN}
\end{center}

The manuscript addresses unified intelligent wireless data and portability
across software implementations. It identifies a failure
mode specific to open, disaggregated RANs: independently implemented
E2SM-KPM reports can be protocol-valid while differing in quantity, scope,
window, clock, counter, missing-value, or preprocessing semantics. KPM-Bridge
places a typed, uncertainty-aware evidence layer between an E2 termination and
portable Near-RT RIC xApps. The design combines a fail-closed semantic-contract
compiler, event-time mapping, split-conformal vector and decision certificates,
and drift-triggered calibration invalidation.

The evaluation uses 201,912 samples from a hash-pinned public ColO-RAN subset,
independent experiment-level splits, seven portable baselines, controlled
implementation profiles, ablations, cluster-bootstrap intervals, and
sensitivity analyses. The manuscript explicitly distinguishes these profiles
from observations of commercial vendors and states the assumptions and limits
of its finite-sample guarantees. Source code, claim-generating outputs, tests,
and reference/claim audits are available in the verified versioned
\href{https://github.com/YassirALKarawi/kpm-bridge-open6g-ran}{KPM-Bridge GitHub repository}.

This work is specific to Open RAN: its contracts bind E2SM-KPM measurement
semantics, its placement is at the E2/Near-RT RIC telemetry boundary, and its
certificate lifecycle is designed for portable xApps. These properties make
the contribution relevant to venues covering open multi-vendor architectures,
intelligent data pipelines, and trustworthy service integration.

Thank you for your consideration.

\closing{Sincerely,}

\end{letter}
\end{document}
